\chapter{Параметры моделей и форматы данных}
\label{app:model-formats}
\ResultPending{Приложение определяет планируемый состав описаний исходных данных и вычислительных опытов. Конкретные файлы, версии моделей и значения параметров будут внесены после выбора данных и разработки метода.}

\section{Описание исходных данных}
Нейродинамика Drosophila и человеческие записи ECAP будут храниться как разные виды данных. Сведения о происхождении должны позволять восстановить, какие наблюдения использованы для настройки, какие --- для проверки, а какие получены расчётным путём. Таблица~\ref{tab:appendix-data-description} задаёт минимальный состав такого описания.

\begingroup\ResearchTableFont
\begin{longtable}{|T{0.25\linewidth}|T{0.67\linewidth}|}
\caption{Планируемый состав описания исходных данных}\label{tab:appendix-data-description}\\\hline
\textbf{Группа сведений} & \textbf{Содержание записи} \\\hline
\endfirsthead
\multicolumn{2}{l}{Продолжение таблицы~\thetable}\\\hline
\textbf{Группа сведений} & \textbf{Содержание записи} \\\hline
\endhead
\endfoot
\endlastfoot
Происхождение & Публикация или репозиторий, версия набора, дата получения, условия использования, контрольная сумма исходного файла. Отдельно отмечаются недоступные сведения. \\\hline
Коннектом & Вид, стадия развития и исследованный экземпляр; анатомический охват, версия реконструкции, идентификаторы нейронов, правила отбора связей и определения их знаков. \\\hline
Функциональные наблюдения & Регистрируемая величина, единицы, частота отсчётов, способ обработки и воздействие. Указывается, относятся ли структура и наблюдения к одному экземпляру либо к совместимым популяциям. \\\hline
Человеческие записи & Обезличенный участник, сеанс и повторное измерение; геометрия контактов, параметры импульса, поза, частота регистрации, способ выделения ответа и измерения его амплитуды. \\\hline
Назначение данных & Настройка модели, внутренняя проверка или независимое испытание. Все записи одного участника связываются общим обезличенным обозначением для контроля разбиения. \\\hline
\end{longtable}
\endgroup

Если наблюдение представляет собой усреднённую флуоресценцию, оно не будет обозначаться как индивидуальная последовательность потенциалов действия без отдельного основания для такого преобразования. Временные шкалы, способы нормировки и преобразования единиц подлежат явному описанию. Совместимость нельзя установить только по совпадению названий областей или типов нейронов.

\section{Описание вычислительного опыта}
Для каждого расчёта планируется сохранять конфигурацию компонентов и связь выходных данных с исходными файлами. Таблица~\ref{tab:appendix-run-description} предназначена для фиксации этих сведений; конкретный программный формат будет выбран при реализации.

\begingroup\ResearchTableFont
\begin{longtable}{|T{0.25\linewidth}|T{0.67\linewidth}|}
\caption{Планируемый состав описания вычислительного опыта}\label{tab:appendix-run-description}\\\hline
\textbf{Группа сведений} & \textbf{Содержание записи} \\\hline
\endfirsthead
\multicolumn{2}{l}{Продолжение таблицы~\thetable}\\\hline
\textbf{Группа сведений} & \textbf{Содержание записи} \\\hline
\endhead
\endfoot
\endlastfoot
Конфигурация & Версия программы, вариант модели, контрольные суммы структуры и параметров, начальное состояние, способ задания случайных чисел и вычислительная точность. \\\hline
Воздействие и время & Форма и амплитуда воздействия, единицы, длительность, шаг интегрирования, частота сохранения, переходный участок и правило его исключения из анализа. \\\hline
Нейродинамический выход & Обозначения переменных, соответствие столбцов нейронам или областям, единицы, временная ось и признак расчётного происхождения. Поведение тела, если оно моделируется, описывается отдельно. \\\hline
Выход человеческой модели & Трансмембранные токи сегментов в амперах, пространственная дискретизация и временная ось; параметры рекрутирования и распространения, геометрия регистрации, напряжение, артефакт, шум и способ измерения амплитуды. \\\hline
Проверка расчёта & Размерности, конечность значений, согласованность времени, повторение с теми же условиями, чувствительность к шагу и статус завершения. Ошибочный расчёт сохраняется как ошибка, а не как допустимое наблюдение. \\\hline
\end{longtable}
\endgroup

Синтетический ряд Drosophila не будет непосредственно обозначаться как ECAP. Целевой электрический сигнал должен быть получен через отдельное описание человеческих волокон и регистрации. До независимой проверки такого описания его выход будет называться ECAP-подобным сигналом. Параметры исходной модели, параметры переносимого блока и параметры целевого оператора будут фиксироваться раздельно.

\chapter{Протокол независимой проверки}
\label{app:independent-validation}
\ResultPending{Ниже приведены требования к будущему протоколу. Окончательные разбиения, численные допуски, необходимое число участников и статистическая процедура будут определены после подтверждения доступа и пригодности данных, до основного сравнения.}

\section{Разделение настройки и испытания}
Проверка будет включать три самостоятельных вопроса: воспроизводит ли исходная модель совместимые функциональные наблюдения; описывает ли человеческий оператор регистрируемый ответ; даёт ли перенос дополнительную ценность относительно контрольных моделей. Положительный ответ на один вопрос не будет подменять проверку остальных. Планируемые условия допуска представлены в таблице~\ref{tab:appendix-validation-record}.

\begingroup\ResearchTableFont
\begin{longtable}{|T{0.25\linewidth}|T{0.67\linewidth}|}
\caption{Планируемые условия независимой проверки}\label{tab:appendix-validation-record}\\\hline
\textbf{Предмет проверки} & \textbf{Сведения, подлежащие фиксации до анализа результата} \\\hline
\endfirsthead
\multicolumn{2}{l}{Продолжение таблицы~\thetable}\\\hline
\textbf{Предмет проверки} & \textbf{Сведения, подлежащие фиксации до анализа результата} \\\hline
\endhead
\endfoot
\endlastfoot
Исходная нейродинамика & Совместимые версии структуры и наблюдений; выбранная регистрируемая величина; данные настройки и отложенные наблюдения. Независимость животных и независимость повторных вычислений описываются раздельно. \\\hline
Физический оператор & Величины и данные калибровки; отдельно --- участники, режимы и записи испытания. Совпадение после подбора параметров не учитывается как независимая проверка на тех же данных. \\\hline
Целевой перенос & Переносимые параметры, адаптеры, одинаковые ограничения вычислений и процедуры выбора моделей. Обязательны модель без переноса и простая модель кривой роста без временного состояния. \\\hline
Контрольные варианты & Раздельные изменения топологии, обученной динамики и источника предобучения; неноцицептивное предобучение; замороженный и дообучаемый перенос. Вариант без физического оператора используется только для диагностики. \\\hline
Независимая оценка & Обезличенные обозначения участников и назначение каждого разбиения; доступная для нового участника информация; основная метрика, практический допуск и способ оценки неопределённости. \\\hline
\end{longtable}
\endgroup

В каждом внешнем разбиении все наблюдения одного участника будут относиться к одной части. Подбор параметров, преобразование временной шкалы, нормировка и калибровка будут выполняться по обучающей части. Если протокол допускает адаптацию на дополнительных данных нового участника, объём этих данных и их отделение от итоговой проверки будут заранее указаны. Такой режим будет оцениваться отдельно от моделирования без адаптации по данным нового участника.

Участников для настройки, самостоятельной проверки оператора и итогового сравнения переноса планируется назначить до анализа результатов. Если проверка оператора используется для выбора модели, обработки или протокола, её участники не будут включены в окончательный тест переноса. Отдельные группы могут быть выделены внутри одного набора до просмотра результатов. Совместное рассмотрение одной отложенной группы допустимо лишь при предварительной фиксации обоих анализов и не будет представлено как независимая репликация. Условия допуска оператора и действия при их невыполнении будут определены до его проверки.

\section{Показатель и интерпретация результата}
Основным показателем предполагается среднеквадратическая ошибка амплитуды кривой роста в микровольтах. Для каждого участника ошибка будет рассчитана по заранее установленным точкам и условиям стимуляции; затем результаты будут усреднены между участниками с одинаковым весом. Повторные записи одного человека не будут увеличивать число независимых участников. Обозначения показателя и условий его применения приведены в разделе~\ref{sec:data-validation}.

Различие методов будет оцениваться на одних и тех же участниках с учётом неопределённости. До сравнения будут определены практически значимый допуск и правило отнесения результата к одному из четырёх исходов: преимущество переноса, подтверждённое отсутствие практически значимого различия, ухудшение либо неопределённость. Одного значения вероятности ошибки статистического критерия недостаточно для утверждения эквивалентности.

При недоступности пригодных человеческих записей результаты будут ограничены вычислительными испытаниями и анализом применимости. При несовместимости исходных наблюдений будет пересмотрена соответствующая заявка на предобучение. При неподтверждённости оператора его выход сохранит статус ECAP-подобного сигнала. Такие обстоятельства не будут представлены как независимая проверка на человеке или как доказательство отсутствия эффекта переноса.
